%% Plantilla de Registro de Decisión de Arquitectura (ADR)
%% Uso: copiar como adr/adr-NNN-<slug>.tex, completar y compilar.
%% Compilar: tectonic -X compile adr-plantilla.tex
\documentclass[11pt,a4paper]{article}
\usepackage{../latex/legasoft}

\lgtitulo{ADR-NNN: \campo{título de la decisión}}
\lgtituloCorto{ADR-NNN}
\lgcodigo{LGS-ADR-NNN}
\lgversion{0.1}
\lgestado{Propuesta}
\lgfecha{\campo{dd-mm-aaaa}}
\lgautor{\lgArquitecto\ (\lgArquitectoRol)}

\begin{document}

% Un ADR es un documento corto: sin portada ni índice.
\begin{center}
{\sffamily\bfseries\LARGE\color{lgPrimario} ADR-\campo{NNN}\par}
\vspace{4pt}
{\sffamily\Large\color{lgPrimario!85!black} \campo{Título de la decisión}\par}
\end{center}
\vspace{8pt}
{\color{lgBorde}\hrule height 0.8pt}
\vspace{10pt}

\begin{tabularx}{\linewidth}{@{}L{2.8cm} Y L{2.4cm} Y@{}}
\sffamily\bfseries\small\color{lgPrimario}Estado & \campo{propuesta / aceptada / rechazada / reemplazada por ADR-NNN} &
\sffamily\bfseries\small\color{lgPrimario}Fecha  & \campo{dd-mm-aaaa} \\
\sffamily\bfseries\small\color{lgPrimario}Autor  & \lgArquitecto &
\sffamily\bfseries\small\color{lgPrimario}Revisores & \lgDirector, \lgQA \\
\end{tabularx}

\vspace{6pt}
{\color{lgBorde}\hrule height 0.4pt}

\begin{porcompletar}[Cómo usar esta plantilla]
Copie este archivo como \ruta{adr/adr-NNN-<slug>.tex} y elimine esta caja. Un
ADR se escribe cuando una decisión condiciona la evolución del producto y
resultaría costosa de revertir. Los ADR no se editan una vez aceptados: si la
decisión cambia, se crea uno nuevo que reemplaza al anterior.
\end{porcompletar}

% =============================================================================
\section*{Contexto}

Situación que obliga a decidir: requisitos afectados, restricciones vigentes y
fuerzas en tensión. Debe poder leerse sin conocer la decisión.

\begin{porcompletar}
\campo{¿Qué problema técnico hay que resolver? ¿Qué requisitos de
\texttt{LGS-REQ-001} lo motivan? ¿Qué restricciones acotan las opciones
disponibles?}
\end{porcompletar}

% =============================================================================
\section*{Alternativas evaluadas}

\begin{center}
\begin{tabularx}{\linewidth}{@{}L{3.4cm} Y Y@{}}
\toprule
\sffamily\bfseries\color{lgPrimario}Alternativa &
\sffamily\bfseries\color{lgPrimario}A favor &
\sffamily\bfseries\color{lgPrimario}En contra \\
\midrule
\campo{Opción A} & \campo{ventajas} & \campo{desventajas} \\
\campo{Opción B} & \campo{ventajas} & \campo{desventajas} \\
\campo{Opción C} & \campo{ventajas} & \campo{desventajas} \\
\bottomrule
\end{tabularx}
\captionof{table}{Alternativas consideradas antes de decidir.}
\end{center}

% =============================================================================
\section*{Decisión}

\begin{nota}[Decisión adoptada]
\campo{Se adopta \ldots\ porque \ldots} La justificación debe conectar la opción
elegida con los requisitos y las restricciones descritos en el contexto, no con
la preferencia del equipo.
\end{nota}

Cada proyecto debe justificar la tecnología que incorpora: si la alternativa
elegida añade complejidad, indique qué beneficio verificable la compensa.

% =============================================================================
\section*{Consecuencias}

\subsection*{Positivas}
\begin{porcompletar}
\campo{Qué se gana: atributos de calidad favorecidos, riesgos mitigados,
capacidades habilitadas.}
\end{porcompletar}

\subsection*{Negativas y compromisos aceptados}
\begin{porcompletar}
\campo{Qué se pierde o se limita, qué deuda técnica se asume y bajo qué
condiciones convendría revisar esta decisión.}
\end{porcompletar}

\subsection*{Impacto en el equipo}
\begin{center}
\begin{tabularx}{\linewidth}{@{}L{4.4cm} Y@{}}
\toprule
\sffamily\bfseries\color{lgPrimario}Rol &
\sffamily\bfseries\color{lgPrimario}Qué cambia para este rol \\
\midrule
\lgDirectorRol   & \campo{impacto en alcance, cronograma o riesgos.} \\
\lgArquitectoRol & \campo{impacto en implementación y mantenimiento.} \\
\lgQARol         & \campo{impacto en la estrategia de pruebas y en la integración.} \\
\bottomrule
\end{tabularx}
\captionof{table}{Consecuencias de la decisión por rol.}
\end{center}

% =============================================================================
\section*{Referencias}

\begin{itemize}
  \item \texttt{LGS-ARQ-001} — Documento de arquitectura del sistema.
  \item \campo{Minutas, pruebas de concepto o documentación externa consultada.}
\end{itemize}

\end{document}
